% =====================================================================
%  Widma EDS — fragment rozdziału 5 (sekcja \ref{sec:widma}).
%  Wstawiany we wspolne/rozdzialy/05_badania.tex po \section{Widma EDS próbek}.
%  Dane: widma-eds/ (README.md, widmo_eds_4kV.xlsx, *_digitalizacja.csv).
%  Rysunki: widma-eds/widmo_eds_4kV.pdf, widma-eds/stabilnosc_zliczen.pdf
% =====================================================================

\subsection{Warunki rejestracji widma}
\label{sec:widma-warunki}

Widmo zarejestrowano 12.06.2026 przy napięciu przyspieszającym
$\UT=\SI{4000}{\volt}$ i~napięciu ekstrakcyjnym $\UB=\SI{1000}{\volt}$.
Impulsy z~detektora promieniowania X przetwarzano cyfrowym procesorem impulsów
KETEK DPP2 obsługiwanym programem ProSpect (v1.1.36); typ detektora nie został
potwierdzony, przy czym zastosowana elektronika współpracuje zwykle z~detektorami
krzemowymi dryfowymi (SDD). Ustawienia toru akwizycji zestawiono w~\tabref{tab:widma-parametry}.

\begin{table}[htbp]\centering
  \caption{Parametry rejestracji widma EDS}\label{tab:widma-parametry}
  \begin{tabular}{@{}ll@{}}\toprule
    Parametr & Wartość \\\midrule
    Napięcie przyspieszające $\UT$ & \SI{4000}{\volt} \\
    Napięcie ekstrakcyjne $\UB$ & \SI{1000}{\volt} \\
    Czas kształtowania (\angielski{peaking time}) & \SI{0.5}{\micro\second} \\
    Liczba kanałów, szerokość kanału & 8192, \SI{10}{\electronvolt} \\
    Czas rzeczywisty / czas życia & \SI{210.7}{\second} / \SI{100.1}{\second} \\
    Częstość zliczeń wejściowa / wyjściowa & ok.~475~kcps / ok.~225~kcps \\
    Czas martwy & \SIrange{52}{54}{\percent} \\\bottomrule
  \end{tabular}
\end{table}

Widmo zbierano w~trybie ciągłym z~kumulacją, a~stan pomiaru zapisano
w~ośmiu chwilach pośrednich. Badaną próbką była płytka aluminiowa, a~pomiar
wykonano z~magnesem ogniskującym umieszczonym na emiterze. Ciśnienie
w~komorze wynosiło ok.~\SI{1e-5}{\hecto\pascal}.
\TODO{prąd emisji $\Ie$ podczas pomiaru, jeśli zapisano}

\subsection{Identyfikacja linii}
\label{sec:widma-linie}

Widmo po \SI{100}{\second} czasu życia przedstawiono na \rysref{fig:widmo-4kV}.
\DOSPRAWDZENIA{Wartości odczytano ze zdjęcia ekranu programu po korekcji
perspektywy, z~dokładnością położenia ok.~\SI{20}{\electronvolt}; przed oddaniem
pracy należy zastąpić je widmem wyeksportowanym z~programu ProSpect.}

\rysunek[0.95\linewidth]{widma-eds/widmo_eds_4kV.pdf}{Widmo EDS zarejestrowane
  przy $\UT=\SI{4}{\kilo\volt}$ i~$\UB=\SI{1}{\kilo\volt}$, czas życia \SI{100}{\second};
  w~powiększeniu zakres pików sumacyjnych \opracowanieWlasne}{fig:widmo-4kV}

W~widmie występują trzy linie charakterystyczne lekkich pierwiastków
(\tabref{tab:widma-linie}). Najsilniejsza jest linia Al~\Ka{} (ok.~$1{,}33\cdot10^{6}$
zliczeń w~kanale \SI{1.49}{\kilo\electronvolt}), a~linie Si~\Ka{} i~Mg~\Ka{} mają
wysokość odpowiednio ok.~\SI{63}{\percent} i~\SI{39}{\percent} piku Al.
Położenia pików zgadzają się z~wartościami tabelarycznymi w~granicach dokładności
odczytu, co potwierdza poprawną kalibrację energetyczną toru akwizycji.

\begin{table}[htbp]\centering
  \caption{Linie zidentyfikowane w~widmie z~\rysref{fig:widmo-4kV}}\label{tab:widma-linie}
  \begin{tabular}{@{}S[table-format=1.2]lS[table-format=1.3]l@{}}\toprule
    {$E$ piku [\si{\kilo\electronvolt}]} & Linia & {$E$ tabel. [\si{\kilo\electronvolt}]} & Pochodzenie \\\midrule
    1.26 & Mg~\Ka & 1.254 & domieszka w~próbce Al \\
    1.48 & Al~\Ka & 1.486 & próbka \\
    1.74 & Si~\Ka & 1.740 & struktura krzemowa\,/\,domieszka w~próbce \\
    2.98 & Al~\Ka{} + Al~\Ka & 2.972 & pik sumacyjny \\
    3.22 & Al~\Ka{} + Si~\Ka & 3.226 & pik sumacyjny \\
    3.46 & Si~\Ka{} + Si~\Ka & 3.480 & pik sumacyjny \\\bottomrule
  \end{tabular}
\end{table}

Ciągłe tło promieniowania hamowania zanika przy energii ok.~\SI{4}{\kilo\electronvolt},
co odpowiada granicy Duane'a--Hunta $E_\mathrm{max}=e\UT$ (\rozdzref{ch:podstawy}).
Elektrony docierają więc do próbki z~pełną energią ok.~\SI{4}{\kilo\electronvolt},
a~napięcie bramki nie zwiększa energii wiązki. Przy tej energii współczynnik
nadnapięcia wynosi ok.~3,1 dla Mg~K, 2,6 dla Al~K i~2,2 dla Si~K, co zapewnia
skuteczne wzbudzenie tych linii; linie K pierwiastków cięższych od wapnia nie
mogą być wzbudzone.

\subsection{Piki sumacyjne i~stabilność zliczeń}
\label{sec:widma-stabilnosc}

Słabe piki przy \SIlist{2.98;3.22;3.46}{\kilo\electronvolt} odpowiadają sumom
energii par fotonów Al i~Si zarejestrowanych w~jednym impulsie. Ich obecność
wynika z~bardzo dużej częstości zliczeń, przy której czas martwy przekracza
\SI{50}{\percent}. W~kolejnych pomiarach należy zmniejszyć prąd wiązki lub
zwiększyć odległość detektora od próbki, tak aby czas martwy nie przekraczał
ok.~\SI{30}{\percent}.

Przebieg częstości zliczeń w~czasie pomiaru przedstawiono na \rysref{fig:widma-stabilnosc}.
Częstości wejściowa i~wyjściowa pozostają stałe w~granicach \SI{\pm2}{\percent},
natomiast średnia częstość zliczeń w~kanale Al~\Ka{} maleje w~ciągu \SI{200}{\second}
z~ok.~$15{,}0\cdot10^{3}$ do $13{,}3\cdot10^{3}$~\si{\per\second}, tj.~o~ok.~\SI{11}{\percent}.
Spadek ten może wynikać ze zmian prądu emisji (podrozdz.~\ref{sec:emisja-magnes}), jednak
przy stałej częstości wejściowej bliskiej nasycenia toru akwizycji nie można
tego rozstrzygnąć bez pomiaru $\Ie$.

\rysunek[0.9\linewidth]{widma-eds/stabilnosc_zliczen.pdf}{Częstość zliczeń wejściowa (ICR),
  wyjściowa (OCR) i~średnia częstość w~kanale Al~\Ka{} w~czasie rejestracji widma
  \opracowanieWlasne}{fig:widma-stabilnosc}

Uzyskane widmo potwierdza, że wiązka elektronów z~emitera CNT wzbudza
w~próbce charakterystyczne promieniowanie X, a~linia Al~\Ka{} wskazana
w~celu pracy jest rejestrowana z~dużą liczbą zliczeń już przy
$\UT=\SI{4}{\kilo\volt}$.
